\section{A conditional energy projection for the weight-matrix operations}\label{sec:energy}
This section is an accounting scenario, not an energy measurement of a training system. It asks how much energy the
weight-matrix operations of a training step would take if the analog reads ran at the efficiency of a published analog-AI
\emph{inference} chip and the updates cost the published or projected pulse energies, and it compares that with an H100.

\paragraph{Assumptions, stated first.}
\begin{enumerate}
\item \textbf{Different read assumptions for accuracy and energy.} The accuracies of Section~\ref{sec:main} are simulated with
      ideal fp32 reads; the energy prices finite-precision analog reads. We have not shown that one operating point gives both
      the reported accuracy and the assumed read efficiency.
\item \textbf{The read efficiency is borrowed from a phase-change inference chip}~\cite{ambrogio2023}: up to 12.4\,TOPS/W
      sustained, measured on one chip, and 6.94\,TOPS/W, the authors' \emph{projected} full-system figure for their five-chip
      system: the chips were not integrated with digital processing, whose time and energy the authors estimate from their
      architecture model, and external components such as system buses and voltage regulators are excluded. We report both. Neither is an ECRAM training array, which also needs
      backward reads, repeated writes, references and Tiki-Taka state.
\item \textbf{The boundary is the weight matrices.} Per weight and per training row: the forward and backward reads (two MACs)
      and the update (device pulses, pulse lines and pulse generation, and the TT family's column reads and digital buffer).
      Attention products, nonlinearities, normalisation, activation storage, data movement, control, reference maintenance and
      static power are not modelled on either side, and the digital exclusions do not cancel the analog ones.
\item \textbf{The GPU side} is H100 board power over its dense BF16 peak~\cite{nvidiah100} at the utilisation of a very large
      training run, 38--43\,\%~\cite{dubey2024}: \RVenGpuPj\,pJ per weight per row (\RVenGpuMeasLo--\RVenGpuMeasHi\,pJ over that
      range; \RVenGpuPlannedLo--\RVenGpuPlannedHi\,pJ over the 30--55\,\% range of our original plan).
\end{enumerate}

\paragraph{Projection.} With the repaired pulse counts (resets counted as applied, the pulse reset's extra reads included), the
GPU-to-analog ratio per weight per row is \RVenChipMin--\RVenChipMax$\times$ at the 12.4\,TOPS/W chip figure and
\RVenSysMin--\RVenSysMax$\times$ at the 6.94\,TOPS/W system figure, for every published device and algorithm. The reads are at
least \RVenChipMvmShareMin\,\% of the analog total: the ratio is set almost entirely by the two borrowed efficiencies, and
changes in the simulated pulse counts barely move it.

\paragraph{Conditional ratios for cells that pass the margin.} A low energy per operation is worth nothing without the
accuracy. PASS means that the whole confidence interval of the gap lies within the stated non-inferiority margin under ideal reads,
not that the losses are equal; we therefore print the ratio only for cells that PASS against the stated baseline
(Table~\ref{tab:rev_energy}), mark FAIL cells as missing the margin and UNRESOLVED cells as ``criterion not established'':
\begin{itemize}
\item against FP-SGD: \RVenEqFpMnistN, \RVenEqFpCifarN{} and \RVenEqFpLmN{} of the \RVenCellsMnistN{} published-device cells per
      task (MNIST, CIFAR, LM);
\item against Adam: \RVenEqAdamMnistN, \RVenEqAdamCifarN{} and \RVenEqAdamLmN{} cells. \textbf{On the LM no analog cell passes
      the margin against Adam, so no conditional ratio exists against the stronger digital baseline.}
\end{itemize}
The comparison is at a fixed training budget; reaching a quality at a different number of steps is not modelled. No joint
finite-read operating point of accuracy and energy, and no energy-to-target training trajectory, is demonstrated; finite
precision, tiling and utilisation, reference maintenance, the digital TT state, update scheduling and full-system costs stay
outside the ratio.

\paragraph{What 100$\times$ would need.} At the central GPU figure, a 100$\times$ ratio needs reads at about
\RVenNeedTopsMin--\RVenNeedTopsMax\,TOPS/W including their periphery, several times the chip figure above.

\paragraph{The memistor.} Widrow's macroscopic memistor moves tens of microjoules of plating charge per level; its ratio is at
most \RVenWidrowRatioMax{} (far worse than the GPU), as expected for wet, centimetre-scale cells.

\begin{table*}[t]\centering\scriptsize\caption{Conditional energy ratios for cells that pass the stated non-inferiority margin under ideal reads: the GPU-to-analog energy ratio per weight per row (12.4 TOPS/W measured chip figure / 6.94 TOPS/W projected system figure of the same inference chip, which includes estimated digital processing), printed only where the cell PASSES against the baseline; fail = the cell FAILS the margin; n.e. = the criterion is not established (UNRESOLVED). No finite-read operating point is demonstrated.}\label{tab:rev_energy}
\begin{tabular}{llcccccc}\toprule & & \multicolumn{2}{c}{MNIST} & \multicolumn{2}{c}{CIFAR} & \multicolumn{2}{c}{LM}\\ Device & Alg. & FP-SGD & Adam & FP-SGD & Adam & FP-SGD & Adam\\\midrule
Li-ECRAM & SGD & fail & fail & fail & fail & fail & fail\\
Li-ECRAM & TTv2 & 32 / 18 & 32 / 18 & fail & fail & n.e. & fail\\
Li-ECRAM & c-TTv2 & 32 / 18 & 32 / 18 & 32 / 18 & 32 / 18 & 32 / 18 & fail\\
Li-ECRAM & AGAD & 32 / 18 & 32 / 18 & 32 / 18 & 32 / 18 & 32 / 18 & fail\\
MO-ECRAM & SGD & fail & fail & fail & fail & fail & fail\\
MO-ECRAM & TTv2 & n.e. & n.e. & fail & fail & fail & fail\\
MO-ECRAM & c-TTv2 & 32 / 18 & 32 / 18 & n.e. & n.e. & fail & fail\\
MO-ECRAM & AGAD & 32 / 18 & 32 / 18 & 32 / 18 & 32 / 18 & fail & fail\\
Protonic-inspired & SGD & fail & fail & fail & fail & fail & fail\\
Protonic-inspired & TTv2 & 32 / 18 & 32 / 18 & fail & fail & n.e. & fail\\
Protonic-inspired & c-TTv2 & 32 / 18 & 32 / 18 & 32 / 18 & n.e. & 32 / 18 & fail\\
Protonic-inspired & AGAD & 32 / 18 & 32 / 18 & 32 / 18 & 32 / 18 & 32 / 18 & fail\\
W/WO$_x$ ECRAM & SGD & fail & fail & fail & fail & fail & fail\\
W/WO$_x$ ECRAM & TTv2 & 32 / 18 & 32 / 18 & fail & fail & fail & fail\\
W/WO$_x$ ECRAM & c-TTv2 & 32 / 18 & 32 / 18 & fail & fail & fail & fail\\
W/WO$_x$ ECRAM & AGAD & 32 / 18 & 32 / 18 & fail & fail & fail & fail\\
\bottomrule\end{tabular}\end{table*}

